\chapter{Methodology and Approach}
\label{cha:MethodologyAndApproach}

This chapter outlines the methodology for developing and benchmarking a cited text span extraction framework. First, it briefly describes the dataset curation process before introducing a hierarchy of text span extraction methods. Lastly, it addresses the necessary evaluation dimensions that will be part of the automatic assessment of the attribution.
\section{Dataset}
In order to curate a dataset of high quality open-domain scientific questions, meeting previously discussed criteria, I will apply an automated pipeline. To guarantee domain alignment of the questions, this pipeline will only rely on papers that are also part of the unarXive 2022 Dataset. Papers will be randomly chosen, to ensure broad coverage and a wide variety of topics across the corpus. Closely mimicking how the process of answering complex scientific questions usually works, with information scattered across numerous sources, the goal of this pipeline is to generate multihop questions, requiring at best complex reasoning capability as well as multiple sources to be reliably and correctly answered. In order to achieve that, not only the chosen paper, but also some of its cited references, will be used as a basis for question generation. Using this method, I aim to produce a dataset with a size of approximately 500 questions, while also collecting additional metadata to later use for further evaluation. 
%Furthermore, a manual verification on a randomly sampled subset will be conducted, to verify that the generated questions satisfy discussed dataset requirements, 
% add manual verification

\section{Cited text span extraction}
Developing a system capable of providing cited text spans alongside its answer introduces an additional constraint in terms of run-time performance, balancing the computation complexity of the retrieval process against potential improvements in retrieved citations. While specific requirements do heavily rely on the domain of application, I want to benchmark a spectrum of different retrieval approaches in terms of correctness and efficiency.
%sowas wie spektrum von dem man wählen kann dunno

Starting with the existing SQuAI extraction method relying on lexical matching, I will further introduce semantic retrieval using Bi-Encoders as well as reasoning capabilities in the form of Cross-Encoders and strategic combinations of those methods to try to further enhance the system's performance. As a gold standard, I will additionally use an LLM, applying its deep reading comprehension and reasoning capabilities.


\section{Span evaluation}
While the main focus of this work is the improvement of the more promising post-hoc attribution of source spans, I will also aim for a comprehensive evaluation of the system as a whole, since post-hoc methods are dependent on the results of prior execution steps. Therefore, further improvements of the span extraction should also be sought there. The evaluation is thus split up into the following parts: 
\begin{enumerate}
\item Evaluating the factual correctness of the generated answers against the dataset's ground truth by applying NLI models and LLM-as-a-judge approaches.
\item Evaluating the retrieval capabilities of the SQuAI system using metrics like precision, recall and F1-Score.
\item Evaluating the performance and latency of extraction variants.
\item Evaluating how suitable the extracted spans are as justification for their respective sentences. This will include using an NLI model to check for logical entailment as well as an LLM-as-a-judge approach measuring metrics like faithfulness, contextual relevance, logical consistency and ambiguity of the attributed text spans.
\end{enumerate}


% 1 - scope to big --> unproblematisch
% 2 - task definition - einen supporting span pro antwortsatz


% 2 - für jede citation korrekten span
%     " find all supporting spans?" --> find one
%     - ein zitat pro satz (limitations) - thema cited text spans nicht multiple citations

% 4- zusätzliches experiment, auf jeden fall die korrekten mit reinfügen und auf 10 auffüllen (testet ob das retrieval wirklich funktioniert) --> schließt limitations durch retrieval aus 
%     - position bias --> random reihenfolge

% 5- - decomposition-based post-hoc attribution (erstmal irgnorieren)
%     - bessere motivation für einzelne varianten

% 6- gut abgedeckt


% 7 - human validation
% 8 - framing als eigenständige arbeit

% 11 - annotation protocol (welche aspekte, wann welcher wert, welche werte möglich) = anweisungen an Annotierer



% benchmarking the dataset
% decontextualized, multihop, (probably all my dimensions)
% does not contain artifacts from the generation itself     -->include
% not only gluing facts from papers together?               -->maybe
% unambiguous?                                              -->maybe